\begin{table}[t]
\tbl{\label{tab:perm:group}}
{\begin{tcolorbox}[tabularx={p{3cm}|p{2cm}|p{6.3cm}|},enhanced,fonttitle=\bfseries\large,fontupper=\normalsize\sffamily,
colback=yellow!10!white,colframe=red!50!black,colbacktitle=blue!30!white,
coltitle=black,title=Creating a Generic Permutation Group]
\hline
Command & Arguments & Purpose \\
\hline
\hline
\verb'-symmetric_group'  &  $n$  &  Symmetric group of degree $n$ and order $n!$.\\
\hline
\verb'-cyclic_group'  &  $n$  & Cyclic group of degree $n$ and order $n$.\\
\hline
\verb'-elementary_abelian_group'  &  $n$  & Elementary abelian group of order $n$ and degree $hp$ where $n=p^h$ and $p$ prime.\\
\hline
\verb'-identity_group'  &  $n$  &  Identity group of degree $n$ and order $1$.\\
\hline
\verb'-dihedral_group'  &  $n$  &  Dihedral group of degree $n$ and order $2n$.\\
\hline
\verb'-bsgs'  &  lab1 lab2 $n$ order base $k$ gens  & Group with a given base and a given set of generators. Here, lab1 and lab2 are labels in ascii and latex, respectively, $k$ is the number of generators, gens is the generators as an Orbiter vector of integers.\\
\hline
\verb'-subgroup_by_generators'  &  lab order $k$ gens  &  Subgroup given by generators. Here, $k$ is the number of generators, gens is an Orbiter vector of integers. \\
\hline
\end{tcolorbox}}
\end{table}
